(1) 远区推迟时间 \(t_r=t-r/c\)，辐射电场横向分量可写为 \(E_\theta(r,\theta,t)=\omega^2p_0\sin\theta\cos(\omega t_r)/(4\pi\varepsilon_0c^2r)\)，正负号取决于 \(\hat{\boldsymbol\theta}\) 的约定；对应磁场 \(B_\phi=E_\theta/c\)。两场横向、同相，\(\boldsymbol E\times\boldsymbol B/\mu_0\) 径向朝外。它们均随 \(1/r\) 衰减，而近场 \(1/r^2\)、\(1/r^3\) 项在远区忽略。

(2) 平均 Poynting 通量 \(\langle S_r\rangle=\varepsilon_0c\langle E_\theta^2\rangle\)，利用 \(\langle\cos^2\rangle=1/2\) 得 \(dP/d\Omega=r^2\langle S_r\rangle=\omega^4p_0^2\sin^2\theta/(32\pi^2\varepsilon_0c^3)\)。轴向 \(\theta=0,\pi\) 无辐射，垂直于偶极轴 \(\theta=\pi/2\) 最强；切勿把 \(4\pi\) 各向同性平均套入角分布。

(3) 用 \(\int\sin^2\theta\,d\Omega=2\pi\int_0^\pi\sin^3\theta\,d\theta=8\pi/3\) 积分，得总平均功率 \(P_{\rm rad}=\omega^4p_0^2/(12\pi\varepsilon_0c^3)\)。若偶极矩振幅与电流联系为 \(p_0=I_0\ell/\omega\)，则 \(P_{\rm rad}=I_0^2\omega^2\ell^2/(12\pi\varepsilon_0c^3)\)。能量随波长变短快速增加是小天线辐射效率的重要原因。

(4) 定义辐射电阻 \(R_{\rm rad}\) 满足平均功率 \(P_{\rm rad}=I_{\rm rms}^2R_{\rm rad}=I_0^2R_{\rm rad}/2\)，故 \(R_{\rm rad}=\omega^2\ell^2/(6\pi\varepsilon_0c^3)=(2\pi/3)Z_0(\ell/\lambda)^2\)，其中 \(Z_0=1/(\varepsilon_0c)\)。低频固定 \(I_0,\ell\) 时 \(R_{\rm rad}\propto\omega^2\to0\)，单位为欧姆。此模型使用沿短偶极均匀电流的赫兹近似，不应与实际三角形电流分布的短偶极天线混同。
